\documentclass[10pt,a4paper]{article}
\usepackage{nexusS5}
\setnexusheader{Dossier enseignant — confidentiel}{S5 — Selim MANSOURI — usage enseignant}
\begin{document}
\nexuscover{SÉANCE 5 --- DOSSIER ENSEIGNANT}{Profil, déroulé, corrigés et barème}{Selim MANSOURI}{Entrée en Troisième}{Mathématiques --- Signes, distributivité et retour à l'unité}
\begin{methodbox}\small\textbf{Programme de référence.} programme de cycle 4 en vigueur pour la classe de Troisième en 2026-2027. Transition visée : acquis de Quatrième 2025-2026 vers la Troisième 2026-2027. \textbf{Attention :} le nouveau programme de cycle 4 publié en 2026 s'applique progressivement : Cinquième en 2026-2027, Quatrième en 2027-2028, Troisième en 2028-2029. Les élèves entrant en Troisième à la rentrée 2026 ne relèvent donc PAS de la nouvelle progression.\end{methodbox}
\begin{keybox}\textbf{DOCUMENT CONFIDENTIEL --- USAGE ENSEIGNANT.} Contient les corrigés, le barème et des données nominatives. Ne pas remettre à l'élève ni le laisser visible pendant la séance.\end{keybox}
\sectiontitle{1. Profil synthétique}
\begin{methodbox}
\textbf{Diagnostic initial.} Diagnostic initial (18 items). Puissances solides ; signes, distributivité, proportionnalité et Pythagore à rectifier ; fractions, équations et trigonométrie à diagnostiquer.
\end{methodbox}
\subsectiontitle{Points d'appui documentés}
\begin{itemize}\item Puissances\end{itemize}
\subsectiontitle{Difficultés documentées}
\begin{itemize}\item Signe d'un produit\item Distributivité\item Retour à l'unité\item Pythagore\item Fractions\item Équations\item Trigonométrie\end{itemize}
\begin{warningbox}\textbf{Limite de la reconstruction.} Les tableaux d'observation séance par séance du dossier individuel sont vierges : il n'existe aucune preuve documentée entre le diagnostic initial et aujourd'hui. Tout statut porté ci-dessous décrit un \emph{état de départ}, jamais une acquisition constatée.\end{warningbox}
\newpage\sectiontitle{2. Matrice de trajectoire --- état avant la séance 5}
\rowcolors{2}{SoftBlue}{white}
\par\noindent\begin{tabularx}{\linewidth}{>{\ttfamily\scriptsize}p{27mm} X >{\scriptsize}p{16mm} >{\scriptsize}p{20mm} >{\scriptsize\centering\arraybackslash}p{9mm}}
\rowcolor{Navy}\color{white}\scriptsize skill\_id & \color{white}Compétence & \color{white}Travaillée en & \color{white}Statut avant S5 & \color{white}Prio. \\
M3E\_PUIS\_01 & Puissances : calcul et produit de puissances de même base & S1 & acquis & OK \\
M3E\_GEO\_01 & Théorème de Pythagore : calcul d'une longueur & S3 & fragile & P1 \\
M3E\_LIT\_01 & $\bullet$ Distributivité : développer k(ax+b) & S2 & fragile & P1 \\
M3E\_LIT\_02 & Réduction d'une expression avec constantes signées & S2 & fragile & P1 \\
M3E\_PROP\_01 & Proportionnalité : quatrième proportionnelle et vitesse & S2, S5 & fragile & P1 \\
M3E\_REL\_01 & $\bullet$ Produit, quotient et priorités avec des nombres relatifs & S1 & fragile & P1 \\
M3E\_EQ\_01 & Équation du premier degré, inconnue dans les deux membres & S2 & non évalué & P2 \\
M3E\_FRAC\_01 & Somme et différence de fractions & S1 & non évalué & P2 \\
M3E\_FRAC\_02 & Produit et quotient de fractions, simplification & S1 & non évalué & P2 \\
M3E\_TRIG\_01 & Choix et emploi d'un rapport trigonométrique (cosinus, sinus) & S4 & non évalué & P2 \\
M3E\_STAT\_01 & Moyenne simple et moyenne pondérée & S5 & non évalué & P3 \\
\end{tabularx}
\par\vspace{1mm}\small\textbf{Lecture de la colonne « Statut avant S5 » :} elle décrit l'état constaté au \emph{positionnement initial}, jamais une acquisition constatée depuis. « acquis » signifie « acquis au positionnement initial » ; « en voie » signifie « en voie d'acquisition ».
\par\small$\bullet$ compétence ciblée par les phases 2 et 3 de la séance. La colonne « Prio. » est \emph{provisoire} : la priorité définitive est produite par \code{tools/analyze_s5.py} après la correction.
\begin{warningbox}\textbf{Effet de récence.} Les compétences suivantes sont retravaillées pendant cette séance, puis évaluées moins d'une heure plus tard : \code{M3E_LIT_01}, \code{M3E_REL_01}, \code{M3E_STAT_01}. Une réussite y mesure une \emph{performance immédiate}, pas une consolidation durable. Le plan de rentrée prévoit pour elles un contrôle différé en semaine 1 ou 2 ; aucun bilan ne doit écrire « acquis durablement » sur cette seule preuve.\end{warningbox}
\sectiontitle{3. Pourquoi ces activités, pour cet élève}
\begin{goalbox}\textbf{Objectif de la séance :} Choisir l'opération à partir de la relation et non des nombres présents — conseil du dossier — en sécurisant d'abord le signe d'un produit puis la distributivité complète.\end{goalbox}
\subsectiontitle{Axe prioritaire (phase 2, 25 min) --- Produit, quotient et priorités avec des nombres relatifs}
Le signe d'un produit figure parmi les priorités du dossier individuel ; c'est le socle de tout le calcul littéral de Troisième.
\par\vspace{1mm}\textbf{Compétence visée :} \code{M3E_REL_01}.
\subsectiontitle{Second axe (phase 3, 20 min) --- Développer avec la distributivité, y compris avec un facteur négatif}
La distributivité figure parmi les priorités du dossier ; l'erreur porte sur la propagation du signe au second terme.
\par\vspace{1mm}\textbf{Compétence visée :} \code{M3E_LIT_01}.
\subsectiontitle{Équilibre commun / individuel}
Sur les 75 minutes : \textbf{51 min} (68\,\%) relèvent des objectifs communs du niveau et \textbf{24 min} (32\,\%) ciblent le profil individuel. Sur l'évaluation : \textbf{14 points} de noyau commun (70.0\,\%) et \textbf{6 points} individualisés (30.0\,\%).
\newpage\sectiontitle{4. Déroulé minute par minute}
\rowcolors{2}{SoftGray}{white}
\par\noindent\begin{tabularx}{\linewidth}{>{\bfseries}p{18mm} p{34mm} X X}
\rowcolor{Navy}\color{white}Temps & \color{white}Phase & \color{white}Conduite & \color{white}Indicateur observable \\
0--10 min & Réactiver les cinq domaines du stage et repérer immédiatemen & Individuel, sans carte de rappel pendant 7 minutes, puis 3 minutes de mise en commun. & Au moins 4 réponses exactes sur 5 ; la certitude est renseignée avant toute vérification. \\
10--35 min & Produit, quotient et priorités avec des nombres relatifs & Rappel bref, exemple guidé au tableau, puis exercices en autonomie. Aide donnée seulement après une tentative écrite. & Trois réussites consécutives sans aide sur la compétence visée. \\
35--55 min & Développer avec la distributivité, y compris avec un facteur & Pas d'exemple guidé : l'élève démarre seul. Intervention uniquement sur demande. & Deux réussites autonomes ; le contrôle est écrit sans qu'on le demande. \\
55--70 min & De la Quatrième à la Troisième : le double produit et l'entr & Travail écrit, correction collective courte à la fin. Ne pas transformer ce temps en cours magistral. & L'élève relie explicitement un acquis de l'année écoulée à un attendu de l'année à venir. \\
70--75 min & Cadrage méthodologique, sans aucune information sur le conte & Lecture des cinq règles, puis remise de l'évaluation. Aucune information sur son contenu. & Les deux engagements sont écrits. \\
75--120 min & Évaluation finale & Individuel, sans document. Partie A environ 10 min, B environ 20 min, C environ 15 min. & Somme des durées cibles : 37 min, marge de 7 min. \\
\end{tabularx}
\newpage\sectiontitle{5. Corrigé du livret de travail}
\subsectiontitle{Phase 1 --- réactivation}
\begin{enumerate}
\item \code{M3E_REL_01} --- $24 + (-15) = \mathbf{9}$. Contrôle : deux facteurs de même signe donnent un produit positif.
\item \code{M3E_FRAC_02} --- $\frac{3\times 8}{4\times 9}=\frac{24}{36}=\mathbf{\frac{2}{3}}$ (ou simplification avant le produit).
\item \code{M3E_LIT_02} --- $8x-6$. Contrôle : $5-8+3+2=2$ et $8\times 1-6=2$.
\item \code{M3E_GEO_01} --- Pythagore : $BC^2=9^2+12^2=81+144=225$, donc $BC=\mathbf{15}$ cm.
\item \code{M3E_STAT_01} --- $\frac{40}{4}=\mathbf{10}$ ; $7 < 10 < 13$ : contrôle concluant.
\end{enumerate}
\subsectiontitle{Phase 2 --- Produit, quotient et priorités avec des nombres relatifs}
\begin{enumerate}
\item $\mathbf{48}$.
\item $4 + (-10) = \mathbf{-6}$.
\item $6\times(-5)=\mathbf{-30}$ : trois facteurs négatifs donnent un produit négatif.
\item $-28 - 4 = \mathbf{-32}$.
\item Le facteur manquant est $\mathbf{-7}$ : le produit étant positif, les deux facteurs sont de même signe.
\item Il a appliqué la règle de l'addition de deux négatifs. Résultat exact : $\mathbf{15}$. Contrôle : $(-5)\times(-3)$ est l'opposé de $(-5)\times 3=-15$.
\end{enumerate}
\minititle{Erreurs probables}
\begin{itemize}
\item \textsc{CONCEPT} --- règle des signes de l'addition appliquée au produit
\item \textsc{METHODE} --- calcul de gauche à droite sans respecter les priorités
\end{itemize}
\subsectiontitle{Phase 3 --- Développer avec la distributivité, y compris avec un facteur négatif}
\begin{enumerate}
\item $\mathbf{15x-20}$.
\item $\mathbf{-6x-21}$.
\item $\mathbf{-20+8x}$.
\item $6x-2+4x+5=\mathbf{10x+3}$.
\end{enumerate}
\minititle{Erreurs probables}
\begin{itemize}
\item \textsc{CONCEPT} --- distributivité appliquée au seul premier terme
\item \textsc{CALCUL} --- signe du second produit erroné
\end{itemize}
\subsectiontitle{Phase 4 --- pont vers Troisième}

\textbf{A.} 1. $(x+3)(x+2)$. \quad 2. $x^2 + 2x + 3x + 6 = x^2 + 5x + 6$.
3. $x=4$ : $(4+3)(4+2)=7\times 6=42$ et $16+20+6=42$. \quad 4. $(x+5)(x+1)=x^2+x+5x+5=x^2+6x+5$.

\textbf{B.} 1. $f(x)=2x-3$. \quad 2. $f(5)=7$ : $7$ est l'\textbf{image} de $5$ ; $f(-2)=-7$.
3. $2x-3=9$ donc $2x=12$ et $x=6$ : $6$ est un \textbf{antécédent} de $9$.
4. Tableau : $-7$ ; $-3$ ; $-1$ ; $3$ ; $7$.

\textbf{Observation à noter :} confondre image et antécédent est une erreur de \textsc{CONCEPT} ;
chercher l'antécédent par essais successifs sans écrire l'équation est une erreur de \textsc{METHODE}.

\newpage\sectiontitle{6. Corrigé de l'évaluation et barème analytique}
\begin{keybox}\textbf{Total : 20 points sur 20.} Noyau commun 14 pts (70.0\,\%), items individualisés 6 pts (30.0\,\%). Somme des durées cibles : 37 minutes.\end{keybox}
\itemhead{A1 --- \code{M3E_REL_01} --- noyau commun}{1}{1}
\begin{graybox}\small\textbf{Réponse attendue :} $7$\end{graybox}
\minititle{Étapes significatives}
\begin{enumerate}\item $(-15)\div 3=-5$\item $(-2)\times(-6)=12$\item $-5+12=7$\end{enumerate}
\minititle{Barème par critère --- la saisie se fait critère par critère}
\par\noindent\begin{tabularx}{\linewidth}{>{\ttfamily\scriptsize}p{9mm} >{\bfseries}p{9mm} >{\ttfamily\scriptsize}p{24mm} >{\scriptsize}p{18mm} X}
\toprule \scriptsize critère & Pts & \scriptsize compétence & \scriptsize preuve & Ce qui est exigé \\ \midrule
c1 & 1 & M3E\_REL\_01 & automatisme & résultat $7$ exact \\
\bottomrule\end{tabularx}
\minititle{Erreurs probables et codes}
\par\noindent\begin{tabularx}{\linewidth}{>{\ttfamily\small}p{26mm} X}
\toprule Code & Description \\ \midrule
CONCEPT & règle des signes non appliquée sur le produit ($-12$) \\
METHODE & opérations effectuées de gauche à droite sans priorité \\
\bottomrule\end{tabularx}
\par\vspace{1mm}\small\textbf{Rôle pour la rentrée :} socle de tout calcul algébrique de Troisième
\par\vspace{2mm}
\itemhead{A2 --- \code{M3E_FRAC_02} --- noyau commun}{1}{1}
\begin{graybox}\small\textbf{Réponse attendue :} $\dfrac{3}{2}$\end{graybox}
\minititle{Étapes significatives}
\begin{enumerate}\item $\frac{4\times 15}{5\times 8}=\frac{60}{40}$\item simplification par $20$ : $\frac{3}{2}$\end{enumerate}
\minititle{Barème par critère --- la saisie se fait critère par critère}
\par\noindent\begin{tabularx}{\linewidth}{>{\ttfamily\scriptsize}p{9mm} >{\bfseries}p{9mm} >{\ttfamily\scriptsize}p{24mm} >{\scriptsize}p{18mm} X}
\toprule \scriptsize critère & Pts & \scriptsize compétence & \scriptsize preuve & Ce qui est exigé \\ \midrule
c1 & 1 & M3E\_FRAC\_02 & automatisme & $\frac{3}{2}$ ou $1{,}5$ exact \\
\bottomrule\end{tabularx}
\minititle{Erreurs probables et codes}
\par\noindent\begin{tabularx}{\linewidth}{>{\ttfamily\small}p{26mm} X}
\toprule Code & Description \\ \midrule
CONCEPT & produit en croix : $\frac{4\times 8}{5\times 15}$ \\
METHODE & réduction au même dénominateur avant un produit \\
\bottomrule\end{tabularx}
\par\vspace{1mm}\small\textbf{Rôle pour la rentrée :} calcul fractionnaire exigé en permanence en Troisième
\par\vspace{2mm}
\itemhead{A3 --- \code{M3E_LIT_02} --- noyau commun}{1}{1}
\begin{graybox}\small\textbf{Réponse attendue :} $4x - 11$\end{graybox}
\minititle{Étapes significatives}
\begin{enumerate}\item $9x-5x=4x$\item $-4-7=-11$\end{enumerate}
\minititle{Barème par critère --- la saisie se fait critère par critère}
\par\noindent\begin{tabularx}{\linewidth}{>{\ttfamily\scriptsize}p{9mm} >{\bfseries}p{9mm} >{\ttfamily\scriptsize}p{24mm} >{\scriptsize}p{18mm} X}
\toprule \scriptsize critère & Pts & \scriptsize compétence & \scriptsize preuve & Ce qui est exigé \\ \midrule
c1 & 1 & M3E\_LIT\_02 & automatisme & $4x-11$ exact \\
\bottomrule\end{tabularx}
\minititle{Erreurs probables et codes}
\par\noindent\begin{tabularx}{\linewidth}{>{\ttfamily\small}p{26mm} X}
\toprule Code & Description \\ \midrule
CONCEPT & $-4-7$ traité comme $-3$ ou $+11$ \\
CONCEPT & $4x$ et $-11$ regroupés \\
\bottomrule\end{tabularx}
\par\vspace{1mm}\small\textbf{Rôle pour la rentrée :} prérequis de la factorisation et des identités remarquables
\par\vspace{2mm}
\itemhead{A4 --- \code{M3E_EQ_01} --- noyau commun}{1}{1.5}
\begin{graybox}\small\textbf{Réponse attendue :} $x = 6$\end{graybox}
\minititle{Étapes significatives}
\begin{enumerate}\item $4x-2x=9+3$\item $2x=12$\item $x=6$\end{enumerate}
\minititle{Barème par critère --- la saisie se fait critère par critère}
\par\noindent\begin{tabularx}{\linewidth}{>{\ttfamily\scriptsize}p{9mm} >{\bfseries}p{9mm} >{\ttfamily\scriptsize}p{24mm} >{\scriptsize}p{18mm} X}
\toprule \scriptsize critère & Pts & \scriptsize compétence & \scriptsize preuve & Ce qui est exigé \\ \midrule
c1 & 1 & M3E\_EQ\_01 & procedure & $x=6$ avec au moins une étape écrite \\
\bottomrule\end{tabularx}
\minititle{Erreurs probables et codes}
\par\noindent\begin{tabularx}{\linewidth}{>{\ttfamily\small}p{26mm} X}
\toprule Code & Description \\ \midrule
CONCEPT & signe non changé lors du passage d'un membre à l'autre \\
CALCUL & $2x=12$ mal résolu alors que le regroupement est correct \\
\bottomrule\end{tabularx}
\par\vspace{1mm}\small\textbf{Rôle pour la rentrée :} prérequis des équations produit et des systèmes
\par\vspace{2mm}
\itemhead{A5 --- \code{M3E_REL_01} --- item individualisé}{1}{1}
\begin{graybox}\small\textbf{Réponse attendue :} $24$\end{graybox}
\minititle{Étapes significatives}
\begin{enumerate}\item $(-7)\times(-4)=28$\item $(-20)\div 5=-4$\item $28-4=24$\end{enumerate}
\minititle{Barème par critère --- la saisie se fait critère par critère}
\par\noindent\begin{tabularx}{\linewidth}{>{\ttfamily\scriptsize}p{9mm} >{\bfseries}p{9mm} >{\ttfamily\scriptsize}p{24mm} >{\scriptsize}p{18mm} X}
\toprule \scriptsize critère & Pts & \scriptsize compétence & \scriptsize preuve & Ce qui est exigé \\ \midrule
c1 & 1 & M3E\_REL\_01 & automatisme & $24$ exact \\
\bottomrule\end{tabularx}
\minititle{Erreurs probables et codes}
\par\noindent\begin{tabularx}{\linewidth}{>{\ttfamily\small}p{26mm} X}
\toprule Code & Description \\ \midrule
CONCEPT & produit de deux négatifs pris négatif \\
METHODE & priorités non respectées \\
\bottomrule\end{tabularx}
\par\vspace{1mm}\small\textbf{Rôle pour la rentrée :} socle de tout le calcul algébrique de Troisième
\par\vspace{2mm}
\itemhead{A6 --- \code{M3E_LIT_01} --- item individualisé}{1}{1}
\begin{graybox}\small\textbf{Réponse attendue :} $-15x + 10$\end{graybox}
\minititle{Étapes significatives}
\begin{enumerate}\item $-5\times 3x = -15x$\item $-5\times(-2)=+10$\end{enumerate}
\minititle{Barème par critère --- la saisie se fait critère par critère}
\par\noindent\begin{tabularx}{\linewidth}{>{\ttfamily\scriptsize}p{9mm} >{\bfseries}p{9mm} >{\ttfamily\scriptsize}p{24mm} >{\scriptsize}p{18mm} X}
\toprule \scriptsize critère & Pts & \scriptsize compétence & \scriptsize preuve & Ce qui est exigé \\ \midrule
c1 & 1 & M3E\_LIT\_01 & automatisme & $-15x+10$ exact \\
\bottomrule\end{tabularx}
\minititle{Erreurs probables et codes}
\par\noindent\begin{tabularx}{\linewidth}{>{\ttfamily\small}p{26mm} X}
\toprule Code & Description \\ \midrule
CONCEPT & second terme non multiplié \\
CALCUL & signe du second produit erroné \\
\bottomrule\end{tabularx}
\par\vspace{1mm}\small\textbf{Rôle pour la rentrée :} prérequis de la double distributivité
\par\vspace{2mm}
\itemhead{B1 --- \code{M3E_GEO_01} --- noyau commun}{2}{4}
\begin{graybox}\small\textbf{Réponse attendue :} $AC = 16$ cm ; $AC < BC$ car l'hypoténuse est le plus grand côté.\end{graybox}
\minititle{Étapes significatives}
\begin{enumerate}\item $BC^2 = AB^2 + AC^2$\item $400 = 144 + AC^2$\item $AC^2 = 256$ donc $AC = 16$\end{enumerate}
\minititle{Barème par critère --- la saisie se fait critère par critère}
\par\noindent\begin{tabularx}{\linewidth}{>{\ttfamily\scriptsize}p{9mm} >{\bfseries}p{9mm} >{\ttfamily\scriptsize}p{24mm} >{\scriptsize}p{18mm} X}
\toprule \scriptsize critère & Pts & \scriptsize compétence & \scriptsize preuve & Ce qui est exigé \\ \midrule
c1 & 1.5 & M3E\_GEO\_01 & justification & relation correcte, théorème cité, $AC=16$ cm \\
c2 & 0.5 & M3E\_GEO\_01 & controle & contrôle de vraisemblance argumenté \\
\bottomrule\end{tabularx}
\minititle{Erreurs probables et codes}
\par\noindent\begin{tabularx}{\linewidth}{>{\ttfamily\small}p{26mm} X}
\toprule Code & Description \\ \midrule
CONCEPT & $20^2 + 12^2$ additionnés : l'hypoténuse n'est pas identifiée \\
JUSTIFICATION & théorème non cité \\
CALCUL & racine carrée de 256 erronée \\
\bottomrule\end{tabularx}
\par\vspace{1mm}\small\textbf{Rôle pour la rentrée :} base de Thalès, de la trigonométrie et du repérage en Troisième
\par\vspace{2mm}
\itemhead{B2 --- \code{M3E_TRIG_01} --- noyau commun}{2}{5.75}
\begin{graybox}\small\textbf{Réponse attendue :} $AB$ : côté adjacent ; $BC$ : hypoténuse ; $\cos\widehat{B}=\frac{5}{13}\approx 0{,}3846$ donc $\widehat{B}\approx 67^\circ$ ; il faudrait le sinus.\end{graybox}
\minititle{Étapes significatives}
\begin{enumerate}\item identification des côtés relativement à l'angle $\widehat{B}$\item $\cos\widehat{B}=\frac{\text{adjacent}}{\text{hypoténuse}}=\frac{5}{13}$\item $\widehat{B}=\arccos(5/13)\approx 67^\circ$\end{enumerate}
\minititle{Barème par critère --- la saisie se fait critère par critère}
\par\noindent\begin{tabularx}{\linewidth}{>{\ttfamily\scriptsize}p{9mm} >{\bfseries}p{9mm} >{\ttfamily\scriptsize}p{24mm} >{\scriptsize}p{18mm} X}
\toprule \scriptsize critère & Pts & \scriptsize compétence & \scriptsize preuve & Ce qui est exigé \\ \midrule
c1 & 1 & M3E\_TRIG\_01 & procedure & côtés correctement nommés et rapport correct \\
c2 & 0.5 & M3E\_TRIG\_01 & automatisme & mesure de l'angle exacte au degré près \\
c3 & 0.5 & M3E\_TRIG\_01 & procedure & sinus identifié à la question 3 \\
\bottomrule\end{tabularx}
\minititle{Erreurs probables et codes}
\par\noindent\begin{tabularx}{\linewidth}{>{\ttfamily\small}p{26mm} X}
\toprule Code & Description \\ \midrule
CONCEPT & cosinus et sinus échangés \\
LECTURE & côté adjacent identifié par rapport au mauvais angle \\
CALCUL & arrondi ou emploi de la calculatrice en radians \\
\bottomrule\end{tabularx}
\par\vspace{1mm}\small\textbf{Rôle pour la rentrée :} trigonométrie réinvestie dès les premiers chapitres de Troisième
\par\vspace{2mm}
\itemhead{B3 --- \code{M3E_PROP_01} --- noyau commun}{2}{4}
\begin{graybox}\small\textbf{Réponse attendue :} $140$ km/h ; moyenne pondérée $=\frac{28+27}{5}=11$, comprise entre $9$ et $14$.\end{graybox}
\minititle{Étapes significatives}
\begin{enumerate}\item $1$ h $30$ min $=1{,}5$ h ; $210 \div 1{,}5 = 140$\item $14\times 2 + 9\times 3 = 55$ ; $55 \div 5 = 11$\end{enumerate}
\minititle{Barème par critère --- la saisie se fait critère par critère}
\par\noindent\begin{tabularx}{\linewidth}{>{\ttfamily\scriptsize}p{9mm} >{\bfseries}p{9mm} >{\ttfamily\scriptsize}p{24mm} >{\scriptsize}p{18mm} X}
\toprule \scriptsize critère & Pts & \scriptsize compétence & \scriptsize preuve & Ce qui est exigé \\ \midrule
c1 & 1 & M3E\_PROP\_01 & procedure & vitesse exacte, conversion de 1 h 30 en 1,5 h visible \\
c2 & 0.7 & M3E\_STAT\_01 & procedure & moyenne pondérée exacte, somme des coefficients correcte \\
c3 & 0.3 & M3E\_STAT\_01 & controle & encadrement entre 9 et 14 écrit \\
\bottomrule\end{tabularx}
\minititle{Erreurs probables et codes}
\par\noindent\begin{tabularx}{\linewidth}{>{\ttfamily\small}p{26mm} X}
\toprule Code & Description \\ \midrule
CONCEPT & $1$ h $30$ converti en $1{,}30$ h \\
METHODE & moyenne pondérée calculée comme une moyenne simple ($11{,}5$) \\
LECTURE & division par 2 au lieu de la somme des coefficients \\
\bottomrule\end{tabularx}
\par\vspace{1mm}\small\textbf{Rôle pour la rentrée :} grandeurs quotients et statistiques, réinvesties toute l'année
\par\vspace{2mm}
\itemhead{B4 --- \code{M3E_PROP_01} --- item individualisé}{2}{4}
\begin{graybox}\small\textbf{Réponse attendue :} $3$ TND l'unité, donc $33$ TND ; $BC = 20$ cm.\end{graybox}
\minititle{Étapes significatives}
\begin{enumerate}\item $21 \div 7 = 3$ puis $3\times 11 = 33$\item $BC^2 = 144 + 256 = 400$ donc $BC = 20$\end{enumerate}
\minititle{Barème par critère --- la saisie se fait critère par critère}
\par\noindent\begin{tabularx}{\linewidth}{>{\ttfamily\scriptsize}p{9mm} >{\bfseries}p{9mm} >{\ttfamily\scriptsize}p{24mm} >{\scriptsize}p{18mm} X}
\toprule \scriptsize critère & Pts & \scriptsize compétence & \scriptsize preuve & Ce qui est exigé \\ \midrule
c1 & 1 & M3E\_PROP\_01 & procedure & prix exact, prix unitaire explicité \\
c2 & 1 & M3E\_GEO\_01 & procedure & $BC = 20$ cm avec la relation écrite \\
\bottomrule\end{tabularx}
\minititle{Erreurs probables et codes}
\par\noindent\begin{tabularx}{\linewidth}{>{\ttfamily\small}p{26mm} X}
\toprule Code & Description \\ \midrule
METHODE & raisonnement additif : $21 + 4$ \\
CONCEPT & hypoténuse non identifiée \\
JUSTIFICATION & prix unitaire non explicité \\
\bottomrule\end{tabularx}
\par\vspace{1mm}\small\textbf{Rôle pour la rentrée :} proportionnalité et Pythagore, outils permanents de la Troisième
\par\vspace{2mm}
\itemhead{C1 --- \code{M3E_GEO_01} --- noyau commun}{4}{9}
\begin{graybox}\small\textbf{Réponse attendue :} $5$ m ; environ $16^\circ$ ; pente $\approx 29{,}2\,\% > 5\,\%$ : non conforme ; $C(5)=545$ TND, $120$ est la part fixe.\end{graybox}
\minititle{Étapes significatives}
\begin{enumerate}\item $4{,}8^2+1{,}4^2=23{,}04+1{,}96=25$ donc longueur $=5$ m\item $\cos\alpha = \frac{4{,}8}{5}=0{,}96$ donc $\alpha\approx 16^\circ$\item $1{,}4 \div 4{,}8 \approx 0{,}2917 = 29{,}2\,\%$, très supérieur à $5\,\%$\item $C(5)=85\times 5+120=425+120=545$ TND ; $120$ ne dépend pas de la longueur\end{enumerate}
\minititle{Barème par critère --- la saisie se fait critère par critère}
\par\noindent\begin{tabularx}{\linewidth}{>{\ttfamily\scriptsize}p{9mm} >{\bfseries}p{9mm} >{\ttfamily\scriptsize}p{24mm} >{\scriptsize}p{18mm} X}
\toprule \scriptsize critère & Pts & \scriptsize compétence & \scriptsize preuve & Ce qui est exigé \\ \midrule
c1 & 1 & M3E\_GEO\_01 & transfert & longueur exacte, théorème cité \\
c2 & 1 & M3E\_TRIG\_01 & transfert & angle exact au degré près, rapport correctement choisi \\
c3 & 1 & M3E\_PROP\_01 & transfert & pente exprimée en pourcentage et conclusion explicite \\
c4 & 1 & M3E\_LIT\_01 & transfert & $C(5)$ exact et rôle du terme constant correctement expliqué \\
\bottomrule\end{tabularx}
\minititle{Erreurs probables et codes}
\par\noindent\begin{tabularx}{\linewidth}{>{\ttfamily\small}p{26mm} X}
\toprule Code & Description \\ \midrule
TRANSFERT & la pente est confondue avec l'angle en degrés \\
CONCEPT & cosinus appliqué avec l'hypoténuse au numérateur \\
JUSTIFICATION & conclusion sur la conformité absente ou non chiffrée \\
CONCEPT & $120$ interprété comme un coût par mètre \\
\bottomrule\end{tabularx}
\par\vspace{1mm}\small\textbf{Rôle pour la rentrée :} tâche complexe type brevet : plusieurs domaines, une décision à justifier
\par\vspace{2mm}
\itemhead{C2 --- \code{M3E_LIT_01} --- item individualisé}{2}{4.25}
\begin{graybox}\small\textbf{Réponse attendue :} $x^2 + 7x + 10$ ; contrôle : $3\times 6 = 18$ et $1 + 7 + 10 = 18$.\end{graybox}
\minititle{Étapes significatives}
\begin{enumerate}\item quatre produits : $x^2 + 5x + 2x + 10$\item réduction : $x^2 + 7x + 10$\end{enumerate}
\minititle{Barème par critère --- la saisie se fait critère par critère}
\par\noindent\begin{tabularx}{\linewidth}{>{\ttfamily\scriptsize}p{9mm} >{\bfseries}p{9mm} >{\ttfamily\scriptsize}p{24mm} >{\scriptsize}p{18mm} X}
\toprule \scriptsize critère & Pts & \scriptsize compétence & \scriptsize preuve & Ce qui est exigé \\ \midrule
c1 & 1 & M3E\_LIT\_01 & procedure & développement complet et réduit \\
c2 & 1 & M3E\_LIT\_01 & controle & contrôle mené sur les deux écritures \\
\bottomrule\end{tabularx}
\minititle{Erreurs probables et codes}
\par\noindent\begin{tabularx}{\linewidth}{>{\ttfamily\small}p{26mm} X}
\toprule Code & Description \\ \midrule
CONCEPT & deux produits seulement \\
CONTROLE & une seule écriture évaluée \\
\bottomrule\end{tabularx}
\par\vspace{1mm}\small\textbf{Rôle pour la rentrée :} double distributivité introduite en phase 4 de cette séance
\par\vspace{2mm}
\newpage\sectiontitle{7. Correspondance diagnostic initial $\leftrightarrow$ évaluation finale}
\rowcolors{2}{SoftBlue}{white}
\par\noindent\begin{tabularx}{\linewidth}{>{\bfseries}p{9mm} >{\ttfamily\scriptsize}p{25mm} >{\ttfamily\scriptsize}p{22mm} >{\ttfamily\scriptsize}p{16mm} X}
\rowcolor{Navy}\color{white}Item & \color{white}\scriptsize skill\_id & \color{white}\scriptsize baseline & \color{white}\scriptsize comparaison & \color{white}Lecture \\
A1 & M3E\_REL\_01 & 3E\_TI\_Q02 & indicative\_skill\_comparison & confrontation indicative seulement : un énoncé parallèle existe au positionnement initial, mais la réponse de l'élève à cet énoncé n'a pas été conservée \\
A2 & M3E\_FRAC\_02 & 3E\_TI\_Q04 & indicative\_skill\_comparison & confrontation indicative seulement : un énoncé parallèle existe au positionnement initial, mais la réponse de l'élève à cet énoncé n'a pas été conservée \\
A3 & M3E\_LIT\_02 & 3E\_TI\_Q08 & indicative\_skill\_comparison & confrontation indicative seulement : un énoncé parallèle existe au positionnement initial, mais la réponse de l'élève à cet énoncé n'a pas été conservée \\
A4 & M3E\_EQ\_01 & 3E\_TI\_Q10 & indicative\_skill\_comparison & confrontation indicative seulement : un énoncé parallèle existe au positionnement initial, mais la réponse de l'élève à cet énoncé n'a pas été conservée \\
A5 & M3E\_REL\_01 & 3E\_TI\_Q02 & indicative\_skill\_comparison & confrontation indicative seulement : un énoncé parallèle existe au positionnement initial, mais la réponse de l'élève à cet énoncé n'a pas été conservée \\
A6 & M3E\_LIT\_01 & 3E\_TI\_Q07 & indicative\_skill\_comparison & confrontation indicative seulement : un énoncé parallèle existe au positionnement initial, mais la réponse de l'élève à cet énoncé n'a pas été conservée \\
B1 & M3E\_GEO\_01 & 3E\_TI\_Q14 & indicative\_skill\_comparison & confrontation indicative seulement : un énoncé parallèle existe au positionnement initial, mais la réponse de l'élève à cet énoncé n'a pas été conservée \\
B2 & M3E\_TRIG\_01 & 3E\_TI\_Q16 & indicative\_skill\_comparison & confrontation indicative seulement : un énoncé parallèle existe au positionnement initial, mais la réponse de l'élève à cet énoncé n'a pas été conservée \\
B3 & M3E\_PROP\_01 & 3E\_TI\_Q18 & indicative\_skill\_comparison & confrontation indicative seulement : un énoncé parallèle existe au positionnement initial, mais la réponse de l'élève à cet énoncé n'a pas été conservée \\
B4 & M3E\_PROP\_01 & 3E\_TI\_Q11 & indicative\_skill\_comparison & confrontation indicative seulement : un énoncé parallèle existe au positionnement initial, mais la réponse de l'élève à cet énoncé n'a pas été conservée \\
C1 & M3E\_GEO\_01 & 3E\_TI\_Q13+3E\_TI\_Q15... & indicative\_skill\_comparison & confrontation indicative seulement : un énoncé parallèle existe au positionnement initial, mais la réponse de l'élève à cet énoncé n'a pas été conservée \\
C2 & M3E\_LIT\_01 & 3E\_TI\_Q07... & not\_comparable & non comparable : contenu introduit pendant la séance 5 elle-même \\
\end{tabularx}
\begin{keybox}\textbf{Aucun écart chiffré ne sera calculé.} Les réponses de l'élève aux 18 questions du positionnement initial n'ont pas été conservées : le dossier n'en garde qu'un statut par domaine. Un statut qualitatif n'est pas une mesure : le convertir en score pour en soustraire un autre produirait un chiffre sans valeur. Le script d'analyse impose donc \code{mastery_delta = null} pour toutes les compétences, et se limite à une confrontation \emph{indicative} entre le statut initial et l'état final.\end{keybox}
\sectiontitle{8. Clés d'interprétation après correction}
\subsectiontitle{Échelle de maîtrise par compétence}
\par\noindent\begin{tabularx}{\linewidth}{>{\bfseries}p{10mm} X}
\toprule Niveau & Signification \\ \midrule
0 & non maîtrisé \\
1 & très fragile \\
2 & partiellement maîtrisé \\
3 & maîtrisé \\
4 & maîtrise solide / transfert réussi \\
\bottomrule\end{tabularx}
\subsectiontitle{Croiser réussite et certitude}
\par\noindent\begin{tabularx}{\linewidth}{>{\bfseries}p{40mm} X}
\toprule Situation & Conduite à tenir \\ \midrule
Juste, certitude 3--4 & acquis disponible : faire transférer sur une situation nouvelle. \\
Juste, certitude 1--2 & presque acquis : faire expliciter la procédure, puis réutiliser. \\
Faux, certitude 1--2 & notion à installer ou procédure à reprendre pas à pas. \\
Faux, certitude 3--4 & priorité absolue : confronter la conviction avant toute reprise technique. \\
\bottomrule\end{tabularx}
\subsectiontitle{Distinguer les niveaux d'erreur}
Une erreur de \textsc{CALCUL} isolée, non reproduite sur les items voisins de même compétence, ne vaut pas non-maîtrise conceptuelle. La chaîne à reconstituer est : \textbf{connaissance} $\rightarrow$ \textbf{choix de méthode} $\rightarrow$ \textbf{exécution} $\rightarrow$ \textbf{justification} $\rightarrow$ \textbf{contrôle}. Les codes d'erreur du barème situent l'écart sur cette chaîne.
\sectiontitle{9. Points d'observation propres à cet élève}
\begin{itemize}\item Indicateurs du dossier : cinq produits relatifs sans erreur, quatre développements complets, trois problèmes proportionnels, trois relations de Pythagore correctes.\item Plusieurs domaines restent à diagnostiquer : recueillir des observations sans conclure à une maîtrise.\item Vérifier que le signe du résultat est annoncé avant le calcul de sa valeur.\end{itemize}
\subsectiontitle{Grille de saisie de la séance}
\par\noindent\begin{tabularx}{\linewidth}{>{\bfseries}p{26mm} X X X}
\toprule Phase & Aide maximale utilisée & Réussite autonome & Observation \\ \midrule
Phase 1 & & & \\[6mm]
Phase 2 & & & \\[6mm]
Phase 3 & & & \\[6mm]
Phase 4 & & & \\[6mm]
\bottomrule\end{tabularx}
\begin{warningbox}\textbf{Avant d'écrire quoi que ce soit sur les acquis de cet élève :} tant que l'évaluation n'est pas corrigée et analysée par \code{tools/analyze_s5.py}, aucune formulation du type « maîtrise désormais », « forte progression » ou « objectif atteint » ne doit être employée. Les formulations admises sont : « à vérifier lors de l'évaluation finale », « observé en séance, à confirmer », « hypothèse de consolidation ».\end{warningbox}
\end{document}
